% ═══════════════════════════════════════════════════════════════
% SPANISCH ARBEITSBLATT - Sequenz 6a: Los lugares
% ═══════════════════════════════════════════════════════════════
% Fachfarbe: #c41e3a (Karminrot)
% Gesamt: 20 Punkte | 4 Aufgaben
% ═══════════════════════════════════════════════════════════════

\documentclass[11pt, a4paper]{article}

% SW-MODUS TOGGLE
\newif\ifswmodus
\swmodusfalse

% PAKETE
\usepackage[utf8]{inputenc}
\usepackage[T1]{fontenc}
\usepackage[ngerman]{babel}
\usepackage{helvet}
\renewcommand{\familydefault}{\sfdefault}

\usepackage[margin=2cm]{geometry}
\usepackage{enumitem}
\usepackage{tabularx}
\usepackage{booktabs}
\usepackage{multirow}
\usepackage{array}

\usepackage{xcolor}
\usepackage{framed}
\usepackage{tcolorbox}
\tcbuselibrary{skins}

\usepackage{fancyhdr}
\usepackage{lastpage}
\usepackage{needspace}

% FARBEN
\definecolor{spanisch-color}{HTML}{c41e3a}
\definecolor{grau-color}{HTML}{888888}
\definecolor{hellgrau-color}{HTML}{f5f5f5}
\definecolor{orange-color}{HTML}{b35c00}

\definecolor{sw-dark}{HTML}{000000}
\definecolor{sw-gray}{HTML}{666666}
\definecolor{sw-white}{HTML}{ffffff}

\ifswmodus
    \colorlet{spanisch}{sw-dark}
    \colorlet{grau}{sw-gray}
    \colorlet{hellgrau}{sw-white}
    \colorlet{orange}{sw-dark}
\else
    \colorlet{spanisch}{spanisch-color}
    \colorlet{grau}{grau-color}
    \colorlet{hellgrau}{hellgrau-color}
    \colorlet{orange}{orange-color}
\fi

\newcommand{\fachfarbe}{spanisch}

% VARIABLEN
\newcommand{\thema}{Los lugares -- Orte in der Stadt}
\newcommand{\sequenz}{06}
\newcommand{\block}{a}
\newcommand{\fach}{Spanisch}
\newcommand{\klasse}{7}
\newcommand{\maxpunkte}{20}
\newcommand{\datum}{\today}

% KOPF- UND FUSSZEILE
\pagestyle{fancy}
\fancyhf{}
\renewcommand{\headrulewidth}{0.4pt}
\renewcommand{\footrulewidth}{0.4pt}

\lhead{\fach{} -- Klasse \klasse}
\chead{AB-\sequenz\block}
\rhead{\datum}

\lfoot{}
\cfoot{}
\rfoot{Seite \thepage\ von \pageref{LastPage}}

% BEFEHLE
\newcommand{\punktebox}[1]{\hfill\fbox{\small \_/#1}}
\newcommand{\luecke}[1]{\rule{#1}{0.4pt}}

\newtcolorbox{merkkasten}[1][]{
    colback=hellgrau,
    colframe=\fachfarbe,
    colbacktitle=white,
    coltitle=\fachfarbe,
    boxrule=1pt,
    arc=0pt,
    left=8pt, right=8pt, top=8pt, bottom=8pt,
    fonttitle=\bfseries,
    attach boxed title to top left={yshift=-2mm, xshift=5mm},
    boxed title style={boxrule=0pt, colback=white},
    title=#1
}

\newtcolorbox{vokabelkasten}{
    colback=white,
    colframe=\fachfarbe,
    boxrule=0.5pt,
    arc=0pt,
    left=5pt, right=5pt, top=5pt, bottom=5pt
}

\newtcolorbox{hinweis}{
    colback=white,
    colframe=orange,
    colbacktitle=white,
    coltitle=orange,
    boxrule=1pt,
    arc=0pt,
    left=5pt, right=5pt, top=5pt, bottom=5pt,
    fonttitle=\bfseries,
    attach boxed title to top left={yshift=-2mm, xshift=5mm},
    boxed title style={boxrule=0pt, colback=white},
    title=Hinweis
}

\newenvironment{aufgabe}[1]{%
    \needspace{5\baselineskip}%
    \nopagebreak[4]%
    \vspace{10pt}
    \noindent\textbf{\textcolor{\fachfarbe}{Aufgabe #1}}
    \nopagebreak[4]%
    \vspace{5pt}
    \nopagebreak[4]%
    \begin{list}{}{\leftmargin=0pt}
    \item[]
}{%
    \end{list}
}

\newcommand{\es}[1]{\textcolor{\fachfarbe}{\textbf{#1}}}

% ═══════════════════════════════════════════════════════════════
% DOKUMENT
% ═══════════════════════════════════════════════════════════════

\begin{document}

% KOPFBEREICH
\begin{center}
    {\Large\bfseries\textcolor{\fachfarbe}{Arbeitsblatt: \thema}}
\end{center}

\vspace{5pt}

\noindent
\begin{tabularx}{\textwidth}{|X|X|X|}
\hline
\textbf{Name:} \luecke{4cm} &
\textbf{Klasse:} \klasse &
\textbf{Datum:} \luecke{2cm} \\
\hline
\end{tabularx}

\vspace{15pt}

% ═══════════════════════════════════════════════════════════════
% MERKKASTEN
% ═══════════════════════════════════════════════════════════════

\begin{merkkasten}[Merkkasten: Orte in der Stadt]
\textbf{Los lugares de la ciudad:}

\vspace{0.5em}

\begin{tabular}{p{5.5cm}p{5.5cm}}
\textbf{el (maskulin)} & \textbf{la (feminin)} \\[0.3em]
\es{el parque} = \luecke{3.5cm} & \es{la plaza} = \luecke{3.5cm} \\[0.8em]
\es{el cine} = \luecke{3.5cm} & \es{la biblioteca} = \luecke{3.5cm} \\[0.8em]
\es{el museo} = \luecke{3.5cm} & \es{la farmacia} = \luecke{3.5cm} \\[0.8em]
\es{el supermercado} = \luecke{3.5cm} & \es{la estación} = \luecke{3.5cm} \\[0.8em]
\es{el hospital} = \luecke{3.5cm} & \\[0.8em]
\es{el centro comercial} = \luecke{3.5cm} & \\
\end{tabular}

\vspace{0.5em}

\textit{Tipp: Wörter auf -o sind oft maskulin (el), auf -a oft feminin (la)!}
\end{merkkasten}

\vspace{10pt}

% ═══════════════════════════════════════════════════════════════
% AUFGABE 1: Zuordnung (5 Punkte)
% ═══════════════════════════════════════════════════════════════

\begin{aufgabe}{1}
Verbinde die spanischen Orte mit der deutschen Bedeutung. \punktebox{5}
\end{aufgabe}

\begin{center}
\renewcommand{\arraystretch}{1.8}
\begin{tabular}{rcp{0.5cm}l}
\es{el parque} & \luecke{2cm} & & A) die Apotheke \\
\es{el cine} & \luecke{2cm} & & B) der Park \\
\es{la farmacia} & \luecke{2cm} & & C) der Bahnhof \\
\es{la estación} & \luecke{2cm} & & D) das Kino \\
\es{el hospital} & \luecke{2cm} & & E) das Krankenhaus \\
\end{tabular}
\end{center}

\vspace{15pt}

% ═══════════════════════════════════════════════════════════════
% AUFGABE 2: Artikel einsetzen (5 Punkte)
% ═══════════════════════════════════════════════════════════════

\begin{aufgabe}{2}
Setze den richtigen Artikel ein: \es{el} oder \es{la}. \punktebox{5}
\end{aufgabe}

\begin{vokabelkasten}
\textit{Artikel:} el (maskulin) | la (feminin)
\end{vokabelkasten}

\vspace{0.5em}

\noindent
\renewcommand{\arraystretch}{1.6}
\begin{tabular}{p{0.5\textwidth}p{0.5\textwidth}}
a) \luecke{1.5cm} museo & b) \luecke{1.5cm} plaza \\
c) \luecke{1.5cm} supermercado & d) \luecke{1.5cm} biblioteca \\
e) \luecke{1.5cm} centro comercial & \\
\end{tabular}

\vspace{15pt}

% ═══════════════════════════════════════════════════════════════
% AUFGABE 3: Lückentext (5 Punkte)
% ═══════════════════════════════════════════════════════════════

\begin{aufgabe}{3}
Was kann man wo machen? Setze den passenden Ort ein. \punktebox{5}
\end{aufgabe}

\begin{vokabelkasten}
\textit{Orte:} el cine | la biblioteca | el supermercado | el hospital | el parque
\end{vokabelkasten}

\vspace{0.5em}

\noindent
\renewcommand{\arraystretch}{1.8}
\begin{tabular}{p{\textwidth}}
a) Ich lese Bücher in \luecke{4cm}. \\
b) Ich kaufe Lebensmittel in \luecke{4cm}. \\
c) Ich schaue Filme in \luecke{4cm}. \\
d) Der Arzt arbeitet in \luecke{4cm}. \\
e) Ich spiele Fußball in \luecke{4cm}. \\
\end{tabular}

\vspace{15pt}

% ═══════════════════════════════════════════════════════════════
% AUFGABE 4: Sätze schreiben (5 Punkte)
% ═══════════════════════════════════════════════════════════════

\begin{aufgabe}{4}
Schreibe 3 kurze Sätze auf Spanisch. Benutze die Orte aus dem Merkkasten. \punktebox{5}
\end{aufgabe}

\begin{hinweis}
Beispiel: \es{En mi ciudad hay un parque.} (In meiner Stadt gibt es einen Park.)
\end{hinweis}

\vspace{0.5em}

\noindent
\renewcommand{\arraystretch}{2.5}
\begin{tabular}{p{\textwidth}}
a) \luecke{13cm} \\
b) \luecke{13cm} \\
c) \luecke{13cm} \\
\end{tabular}

% ═══════════════════════════════════════════════════════════════
% PUNKTEÜBERSICHT
% ═══════════════════════════════════════════════════════════════

\vspace{20pt}
\hrule
\vspace{10pt}

\begin{flushright}
\textbf{Punkte:} \luecke{1cm} / \maxpunkte
\end{flushright}

\end{document}
